\section[GER]{Bundestagswahl 2017}

{
\usebackgroundtemplate{
  \tikz[overlay,remember picture] 
  \node[opacity=0.8, at=(current page.south east),anchor=south east,inner sep=0pt] {
    \includegraphics[height=0.9\textheight]{assets/Germany.png}};
}
\begin{frame}{Bundestagswahl 2017}\framesubtitle{German Election}\titlemark{Picture by Dimitri Averni, CC BY-SA 3.0}
 \begin{itemize}
    \item 88499 Wahlbezirke from 299 Wahlkreise 
    \item 6 features (voting results of major parties)
    \item 16 federal states as labels
\end{itemize}
\end{frame}
}

\subsection{Embedding}

\begin{frame}{Wasserstein t-SNE Embedding}
\LargeFigure{GER/Embedding.pdf}
\end{frame}

\begin{frame}{}
\begin{figure}[!ht]
  {\includegraphics[width=.95\textwidth]{GER/FeatureMeans.pdf}}\\
  {\includegraphics[width=.95\textwidth]{GER/FeatureStds.pdf}}
\end{figure}
\end{frame}

\subsection{Analysis}


\begin{frame}{Covariance Analysis}
\LargeFigure{GER/Covariances.pdf}
\end{frame}


\begin{frame}{Correlation Analysis}
\LargeFigure{GER/Correlation.pdf}
\end{frame}


\subsection{Transformation}

\begin{frame}{Variance Stabilization}\framesubtitle{Data means correlate with data variance}
		\begin{center}
		\includegraphics[height=0.5\textheight]{GER/MeanStdCorr.pdf}
		\end{center}
        \begin{proposition}
        If $Y\sim B(n, p)$ (Binominal distribution), then $T(Y) = \sqrt{n} \arcsin(\sqrt{\frac{Y}{n}})$ is a variance stabilizing transformation with $\var(T(Y)) \approx \frac{1}{4}$
        \end{proposition}
\end{frame}


\begin{frame}{Embedding after transformation}
\LargeFigure{GER/Transformation.pdf}
\end{frame}